% 第11回　プレゼン資料作成（記入例・模範解答）
\session{11}{2}{12月9日（水）}{プレゼン資料作成}{AIと構成編集}{yes}

\goals{
  \item 「課題、顧客、コンセプト、空間、体験、効果、実現の道筋」の流れで提案を組み立てられる
  \item 1枚1メッセージのスライドを作れる
  \item AIが出したコピー案から選び、自分たちの言葉に直せる
  \item AIの活用と自分たちの思考の反映を説明するスライドを用意できる
}

\afillin{発表時間　\underline{\makebox[12mm]{\ans{10}}} 分　／　聞き手}{企業の経営企画と店舗運営の担当者}

\work[120mm]{ワーク1}{構成案と、各スライドのメッセージ}
\begin{promptbox}[構成を組み立てる]
企業担当者向けに店舗設計を提案する（発表時間）分のプレゼンの構成を作ってください。聞き手は経営企画と店舗運営の担当者です。各スライドについて、「スライドの見出し（1行のメッセージ）」と「載せる内容」を示してください。提案の要約：（改訂版の要約）
\end{promptbox}
\begin{wstabh}{colspec={Q[c,wd=6mm,m] Q[l,wd=17mm,m] X[3,l,m] X[2,l,m] Q[l,wd=14mm,m]},row{2-Z}={ht=10mm}}
  & 流れ & 見出し（1行のメッセージ） & 載せる内容・使う素材 & 担当 \\
  1 & 課題 & \af{若い世代は、車ではなく「店」から離れている} & \af{課題の整理、アンケートの結果} & \af{自分} \\
  2 & 顧客 & \af{2人の来店客が、店の前でためらっている} & \af{ペルソナカード（第6回）} & \af{Ｂ} \\
  3 & コンセプト & \af{見るだけ、大歓迎。} & \af{コンセプト文、誰に・何を・どのように} & \af{Ｄ} \\
  4 & 空間 & \af{店の真ん中に、子どもの居場所をつくる} & \af{ゾーニング図、生成イメージ（第8回）} & \af{Ｃ} \\
  5 & 体験 & \af{声をかけるのは、お客さまが呼んだとき} & \af{ジャーニーマップ、改善前後の図（第9回）} & \af{Ｃ} \\
  6 & 効果 & \af{来店のハードルを下げ、相談の質を上げる} & \af{根拠（アンケート、事例）} & \af{Ｂ} \\
  7 & 実現の道筋 & \af{明日からできることから、3段階で} & \af{段階ごとの施策と費用の大小} & \af{Ｄ} \\
  8 & AI活用 & \af{AIは発想を広げ、判断は私たちがした} & \af{AIを使った場面、採用・不採用の例} & \af{自分} \\
  9 & まとめ & \af{「行ってみたい」と思える店へ} & \af{コンセプトの再掲、質疑へ} & \af{Ｄ} \\
  10 & & & & \\
\end{wstabh}

\work[60mm]{ワーク3}{キャッチコピー（20案から3つ、そして1つ）}
\begin{promptbox}[コピーの案を出す]
次のコンセプトを表すキャッチコピーを20案出してください。10字以内のものと20字以内のものを半分ずつ。企業の担当者に向けた落ち着いた表現と、来店客に向けた親しみやすい表現の両方を含めてください。コンセプト：（コンセプト文）
\end{promptbox}
\begin{wstabh}{colspec={X[2,l,m] X[3,l,m]},row{2-Z}={ht=10mm}}
  残した3案 & 選んだ理由・気になる点 \\
  \af{見るだけ、大歓迎。} & \af{ペルソナBの不安に直接答える。短く覚えやすい} \\
  \af{車選びは、公園みたいに。} & \af{コンセプトの「公園ラウンジ」が伝わる。少しわかりにくい} \\
  \af{あなたのペースで、決めていい。} & \af{企業担当者にも伝わる落ち着いた表現。やや平凡} \\
\end{wstabh}
\abox[決定したコピー（自分たちの言葉に直した点も）]{18mm}{「見るだけ、大歓迎。」　AI案は「見るだけでも大歓迎です」。言い切りにして短くし、句点で余韻を残した。}

\work[70mm]{AI活用スライド}{「AIをどう活用し、自分たちの思考をどう反映させたか」の下書き}
\begin{wstab}{colspec={Q[l,wd=34mm,m] X[l,m]},row{1-Z}={ht=15mm}}
  \hc{AIを使った場面} & \af{ペルソナ案、コンセプトの発想（30案）、空間の言語化と画像生成、店内行動の予測、企業側の視点での批判、構成とコピー案} \\
  \hc{採用したもの} & \af{「呼ぶまで話しかけない」という反転の発想、つまずきの場面、企業側の指摘6つのうち4つへの対応} \\
  \hc{捨てたもの\newline とその理由} & \af{課題の対象外のペルソナ、出典のない来店増の数値、実在メーカーのロゴに似た生成画像} \\
  \hc{自分たちが\newline 判断したこと} & \af{どのペルソナに絞るか、つまずきの優先順位、どの指摘を割り切るか。判断の基準は常に「ペルソナの言葉」と「企業の課題」} \\
\end{wstab}

\work[50mm]{リハーサル}{相互フィードバック}
\begin{wstabh}{colspec={X[l,m] X[l,m] Q[c,wd=18mm,m]},row{2-Z}={ht=16mm}}
  良かった点 & 改善点 & 所要時間 \\
  \af{見出しだけで話の流れがわかる。ペルソナの言葉で始まるので引き込まれる} & \af{空間のスライドの文字が多い。図を大きくして説明は口頭で} & \ans{12}\,分 \\
  \af{改善前後の図で、変化がひと目でわかる} & \af{時間超過。効果と実現の道筋を1枚にまとめる} & \ans{9.5}\,分 \\
\end{wstabh}

\begin{nexttime}
  \item スライドを完成させる。「AIをどう活用し、自分たちの思考をどう反映させたか」のスライドを必ず含める
  \item 発表時間内に収まるよう、通しで練習する
\end{nexttime}
\aimemoA{プレゼンの構成案、キャッチコピー案}
{「経営企画と店舗運営の担当者向けに、10分の構成を。各スライドの見出しと内容を」「コピーを20案、10字以内と20字以内を半分ずつ」}
{構成案の流れは参考にしたが、見出しはすべて自分たちで書き直した。コピーは20案から3案に絞り、表現を直して決めた。}
{スライドの数値と事例の出典を、それぞれの一次資料で確認した。生成画像に「AIによる生成イメージ」と入れた}

\begin{point}
  \item 見出しを読むだけで提案の流れがわかるか（1枚1メッセージ）を確認する。見出しが「空間」「体験」のような単語だけなら、文にさせる。
  \item 流れが「課題→顧客→コンセプト→空間→体験→効果→実現の道筋」になっているか。効果の根拠と実現の道筋が抜けやすい。
  \item AI活用スライドは、使った場面の列挙で終わらせず、「捨てたもの」と「自分たちの判断」が書けているかを評価する。発表会の評価の観点にも入っている。
  \item リハーサルでは時間を必ず計らせる。超過したグループには、削るスライドを決めさせる。
\end{point}
